\documentclass{beamer}

% Use Feather theme
\usetheme{feather}
\usecolortheme{default}

% Packages
\usepackage[utf8]{inputenc}
\usepackage{amsmath}
\usepackage{amssymb}
\usepackage{booktabs}
\usepackage{xcolor}
\usepackage{tikz}

% Theme customization
\setbeamerfont{title}{size=\Large}
\setbeamerfont{subtitle}{size=\small}

% Remove background graphics - use plain white background for all slides
\setbeamertemplate{background}{%
  \begin{tikzpicture}[remember picture, overlay]
    \fill[white] (current page.south west) rectangle (current page.north east);
  \end{tikzpicture}
}

% Information
\title{iRDRC: Markov Decision Process}
\subtitle{Deep Dive}
\author{Joel James \& Majida Bee}
\date{08 September 2026}
\institute{IIIT Kottayam --- BTP Discussion 2}

\begin{document}

% ============================================================================
% SLIDE 1: TITLE SLIDE
% ============================================================================
\begin{frame}[plain]
  \titlepage
  \vspace{-0.5cm}
  \begin{center}
    \small\textit{Focus: Complete mathematical formulation of the MDP, why it works, and how to implement it.}
  \end{center}
\end{frame}

% ============================================================================
% SLIDE 2: RECAP FROM DISCUSSION 1
% ============================================================================

\begin{frame}{Recap: The iRDRC Problem}
  \begin{itemize}
    \item \textbf{Problem:} AV needs radar + communication on shared resource
    \item \textbf{Environment:} Weather, road, speed, moving objects (binary states)
    \item \textbf{Objective:} Maximize throughput while minimizing miss-detection risk
    \item \textbf{Approach:} Adaptive learned policy via Deep Q-Network (DQN)
    \item \textbf{Key Insight:} Fixed scheduling fails under dynamic uncertainty
  \end{itemize}

  \vspace{0.3cm}
  \begin{center}
    \small
    \textbf{Meeting 1:} System model, environment model (Eq. 1), MDP intro\\
    \textbf{Today:} Deep understanding of MDP, why each component matters, DQN algorithm
  \end{center}
\end{frame}

% ============================================================================
% SLIDE 3: WHY MARKOV DECISION PROCESS?
% ============================================================================
\begin{frame}{Why Markov Decision Process (MDP)?}
  \textbf{The Markov Property:} The future depends only on the current state, not on history.

  \vspace{0.2cm}
  \textbf{Why this matters for iRDRC:}

  \begin{enumerate}
    \item Queue state $d$ captures current backlog; only $d$ matters, not arrival history
    \item Channel state $c$ is independent each step; history doesn't predict next $c$
    \item Environmental factors $(r, w, v, m)$ are resampled; each step is a fresh draw
    \item Radar detection is perfect; if on and event exists, we catch it
  \end{enumerate}

  \vspace{0.3cm}
  \textbf{Result:} Given current state $s$ and action $a$, the next state $s'$ is probabilistically determined. This allows us to use RL algorithms that exploit the Markov property.
\end{frame}

% ============================================================================
% SLIDE 4: THE MDP TUPLE
% ============================================================================
\begin{frame}{The MDP Tuple: $\langle S, A, P, R, \gamma \rangle$}
  A Markov Decision Process is defined by five components that describe a complete sequential decision problem.

  \vspace{0.1cm}
  \scriptsize
  \begin{tabular}{|p{1.6cm}|p{1.4cm}|p{3.2cm}|p{2.8cm}|}
    \hline
    \textbf{Component} & \textbf{Symbol} & \textbf{Meaning (iRDRC)} & \textbf{Example} \\
    \hline
    State Space & $S$ & All possible observations & $(d=5, c=0, r=0, w=1, v=1, m=0)$ \\
    \hline
    Action Space & $A$ & All possible choices & $\{0: \text{comm}, 1: \text{radar}\}$ \\
    \hline
    Transition & $P(s'|s,a)$ & How state evolves & $P(d=4 | d=5, a=0, c=0)$ \\
    \hline
    Reward & $R(s,a,s')$ & Immediate reward & $+2$ (good comm), $-50$ (missed event) \\
    \hline
    Discount & $\gamma$ & Weight on future rewards & $\gamma \approx 0.99$ \\
    \hline
  \end{tabular}

  \normalsize
\end{frame}

% ============================================================================
% SLIDE 5: STATE SPACE (EQUATION 2)
% ============================================================================
\begin{frame}{State Space (Equation 2)}
  \textbf{Definition:}
  $$S = \{ (d, c, r, w, v, m) : d \in \{0,\ldots,10\}, c,r,w,v,m \in \{0,1\} \}$$

  \vspace{0.2cm}
  \textbf{State Space Size Calculation:}
  \begin{itemize}
    \item $d$ (queue): 11 possible values $(0, 1, 2, \ldots, 10)$
    \item $c$ (channel): 2 possible values $(0 = \text{good}, 1 = \text{bad})$
    \item $r$ (road): 2 possible values
    \item $w$ (weather): 2 possible values
    \item $v$ (speed): 2 possible values
    \item $m$ (moving objects): 2 possible values
  \end{itemize}

  \vspace{0.2cm}
  \textbf{Total: $11 \times 2^5 = 352$ states}

  \vspace{0.15cm}
  \small 352 states is manageable but large enough to benefit from function approximation (neural networks) instead of tabular methods.
\end{frame}

% ============================================================================
% SLIDE 6: ACTION SPACE
% ============================================================================
\begin{frame}{Action Space (Simple but Powerful)}
  \textbf{Definition:}
  $$A = \{0, 1\}$$

  \vspace{0.2cm}
  \scriptsize
  \begin{center}
    \begin{tabular}{|p{1.2cm}|p{1.8cm}|p{4.2cm}|}
      \hline
      \textbf{Action} & \textbf{Mode} & \textbf{What Happens} \\
      \hline
      $a = 0$ & Communication & Transmit queued packets; channel determines count \\
      \hline
      $a = 1$ & Radar & Scan for obstacles; detect if event exists \\
      \hline
    \end{tabular}
  \end{center}

  \vspace{0.15cm}
  \normalsize
  \textbf{Performance Impact:}
  \small
  \begin{itemize}
    \setlength{\itemsep}{0.05cm}
    \item Communication: Increases throughput, increases miss-detection risk
    \item Radar: Decreases throughput, decreases miss-detection risk
  \end{itemize}

  \vspace{0.1cm}
  \small\textbf{Central Design Decision:} Each time step, the AV chooses exactly one mode. This binary action is the lever the learning algorithm will adjust.

  \normalsize
\end{frame}

% ============================================================================
% SLIDE 7: REWARD FUNCTION (EQUATION 3 / EQUATION 4)
% ============================================================================
\begin{frame}{Reward Function (Correct Definition)}
  \textbf{The Learning Signal:} Translates engineering objectives (throughput + safety) into numbers a learning algorithm can optimize.

  \vspace{0.1cm}
  \scriptsize
  \begin{tabular}{|p{1.2cm}|p{1.4cm}|p{1.6cm}|p{1.0cm}|p{2.6cm}|}
    \hline
    \textbf{Action} & \textbf{Channel} & \textbf{Event} & \textbf{Reward} & \textbf{Interpretation} \\
    \hline
    $a=0$ (Comm) & $c=0$ (good) & $\ominus$ (NO) & $+r_1 = +2$ & Success: transmit $\nu_1=4$ packets \\
    \hline
    $a=0$ (Comm) & $c=1$ (bad) & $\ominus$ (NO) & $+r_2 = +1$ & Limited: transmit $\nu_2=2$ packets \\
    \hline
    $a=0$ (Comm) & any & $\oplus$ (YES) & $-r_3 = -50$ & Catastrophic: missed obstacle \\
    \hline
    $a=1$ (Radar) & any & $\oplus$ (YES), $b$ unfav. & $+r_4(b+1) = +5(b+1)$ & Success: reward scales with risk \\
    \hline
    $a=1$ (Radar) & any & $\ominus$ (NO) & $0$ & Defensive: no threat encountered \\
    \hline
  \end{tabular}

  \normalsize
  \vspace{0.1cm}
  \small\textbf{Key Design Insights:}
  \begin{itemize}
    \setlength{\itemsep}{0.05cm}
    \item \textbf{Rewards ONLY trigger when NO event} for communication mode ($+r_1, +r_2$ for $\ominus$)
    \item \textbf{Penalty applied} when AV is in communication mode AND event occurs ($-r_3$ regardless of channel)
    \item \textbf{Radar reward scales} with $b$ = number of unfavorable conditions detected (max $+25$ when all 4 factors unfavorable)
    \item $|-50| \gg |r_1|, |r_2|$: Missing event is 25--50× costlier than communication loss
  \end{itemize}

  \normalsize
\end{frame}

% ============================================================================
% SLIDE 8: OBJECTIVE AND DISCOUNTED RETURN
% ============================================================================
\begin{frame}{Objective: Discounted Return (Equation 4)}
  \textbf{Goal:} Find policy $\pi$ that maximizes expected cumulative reward over many time steps.

  \vspace{0.2cm}
  \textbf{Discounted Return (Value of Policy $\pi$):}
  $$G(\pi) = E\left\{ \sum_{t=0}^{T} \gamma^t r_{t+1}(\pi) \right\}$$

  \vspace{0.2cm}
  \textbf{Components:}
  \begin{itemize}
    \item $r_{t+1}(\pi)$: Immediate reward at time $t+1$ under policy $\pi$
    \item $\gamma \in (0,1)$: Discount factor (paper doesn't specify; typical $\approx 0.99$)
    \item $\gamma^t$: Exponential decay; rewards 10 steps away worth $\gamma^{10} \approx 0.90\times$ immediate
    \item $T$: Time horizon (episode length; likely 500--1000 steps)
  \end{itemize}

  \vspace{0.2cm}
  \textbf{Optimal Policy:}
  $$\pi^* = \arg\max_\pi G(\pi)$$

  \small The policy that maximizes long-term payoff. If communication now leads to a crash later, the $-50$ penalty far outweighs the $+1$ or $+2$ gain.
\end{frame}

% ============================================================================
% SLIDE 9: DISCOUNT FACTOR DEEP DIVE
% ============================================================================
\begin{frame}{Why Discount the Future? ($\gamma$ Explained)}
  \textbf{What $\gamma$ Controls:}

  \begin{itemize}
    \item \textbf{$\gamma \to 1.0$ (e.g., 0.99):} Future rewards nearly as valuable as immediate rewards. Policy is far-sighted. Over $T=1000$, reward at step 999 is still worth $\gamma^{999} \approx 0.0005$ of immediate reward.

    \item \textbf{$\gamma \to 0.0$ (e.g., 0.5):} Only immediate rewards matter. Policy is myopic; will choose any immediate gain even if it ruins future prospects.
  \end{itemize}

  \vspace{0.2cm}
  \textbf{For iRDRC:}
  \begin{itemize}
    \item $\gamma \approx 0.99$ makes sense: a crash today affects all future steps
    \item If $\gamma = 0.5$, policy might communicate recklessly to gain $+2$, knowing $-50$ crash is only 50\% as important
    \item $\gamma \approx 0.99$ ensures policy balances immediate throughput against long-term safety
  \end{itemize}
\end{frame}

% ============================================================================
% SLIDE 10: ACTION-VALUE FUNCTION & POLICY
% ============================================================================
\begin{frame}{From MDP to Policy: Q-Values and $\pi$}
  \footnotesize
  \textbf{Action-Value Function $Q(s, a)$:} Expected discounted return from state $s$, action $a$, following optimal policy.

  \vspace{0.08cm}
  $$Q^*(s, a) = E\left\{ r + \gamma \max_{a'} Q^*(s', a') \right\}$$

  \vspace{0.08cm}
  \textbf{Intuition:}
  \begin{itemize}
    \setlength{\itemsep}{0.03cm}
    \item $Q(d=0, a=0)$ high: queue empty, communication valuable
    \item $Q(d=10, a=0)$ lower: queue full
    \item $Q(w=1, v=1, a=0)$ very low: unfavorable weather/speed, communication
    \item $Q(w=1, v=1, a=1)$ higher: radar protects from risk
  \end{itemize}

  \vspace{0.08cm}
  \textbf{Optimal Policy $\pi^*$:} $\pi^*(s) = \arg\max_{a \in A} Q^*(s, a)$

  \scriptsize Pick action with max Q-value. \textbf{Challenge:} Learn $Q^*$ from experience.

  \normalsize
\end{frame}

% ============================================================================
% SLIDE 11: STATE TRANSITIONS
% ============================================================================
\begin{frame}{Environment Dynamics: State Transitions $P(s'|s,a)$}
  \footnotesize
  \textbf{Queue Dynamics (Deterministic):}
  $$d'_{t+1} = \min\left( d_t - \text{transmission} + \text{arrivals}, D \right)$$
  Transmission: $\nu_1=4$ if $c=0$ (good), $\nu_2=2$ if $c=1$ (bad). Arrivals $\sim \text{Poisson}(\lambda_d=1)$.

  \vspace{0.08cm}
  \textbf{Channel Dynamics (Stochastic, Independent):}
  $$c'_{t+1} \sim \text{Bernoulli}(p_c=0.1)$$
  $P(c'=1 | c, a) = 0.1$ (independent of $c$ and $a$).

  \vspace{0.08cm}
  \textbf{Environmental Factors (Stochastic, Independent):}
  Each of $r, w, v, m$ is resampled: $P(x'=1) = 1 - \tau_x$. All factors independent (Markov property).

  \vspace{0.08cm}
  \textbf{Unexpected Event Generation (Stochastic):}
  $$\oplus_t \sim \text{Bernoulli}(P(\oplus))$$
  \scriptsize If $a_t=1$ (radar) and $\oplus_t=1$, detect. If $a_t=0$ (communication) and $\oplus_t=1$, miss ($-50$ penalty).

  \normalsize
\end{frame}

% ============================================================================
% SLIDE 12: TABULAR Q-LEARNING VS DQN
% ============================================================================
\begin{frame}{Tabular Q-Learning vs. Deep Q-Network (DQN)}
  \scriptsize
  \begin{tabular}{|p{1.2cm}|p{3.3cm}|p{3.3cm}|}
    \hline
    \textbf{Aspect} & \textbf{Tabular Q-Learning} & \textbf{Deep Q-Network} \\
    \hline
    Storage & 352 table entries & Neural network weights \\
    \hline
    Update & $Q \gets Q + \alpha[r + \gamma \max Q(s',a') - Q]$ & Same rule, network output \\
    \hline
    Converge & $\sim$280 episodes & $\sim$170 episodes \\
    \hline
    Scalability & 352 states OK; larger impractical & Scales to large state spaces \\
    \hline
    Data Efficiency & One update per experience & Mini-batch updates (variance reduction) \\
    \hline
  \end{tabular}

  \vspace{0.1cm}
  \small\textbf{Bottom Line:} For iRDRC with 352 states, both work. DQN is faster. Both exploit MDP structure; learning from experience is enough.

  \normalsize
\end{frame}

% ============================================================================
% SLIDE 13: DQN NETWORK UPDATE
% ============================================================================
\begin{frame}{DQN Network Update (Equation 5)}
  \textbf{Update Rule:}
  $$\theta_{t+1} = \theta_t + \alpha \left[ y_t - Q(s_t, a_t; \theta_t) \right] \nabla Q(s_t, a_t, \theta_t)$$

  \vspace{0.2cm}
  \textbf{Where:}
  \begin{itemize}
    \item $\theta_t$: Network weights at time $t$
    \item $\alpha$: Learning rate (e.g., 0.001)
    \item $y_t = r_t + \gamma \max_{a'} Q(s_{t+1}, a'; \theta^-_t)$: Target value
    \item $\theta^-_t$: \textbf{Target network} weights (copied from online network periodically, e.g., every 10k steps)
    \item $\nabla Q(\cdots)$: Gradient of Q-value w.r.t. weights
  \end{itemize}

  \vspace{0.2cm}
  \textbf{Key Innovation: Target Network}

  Why two networks? Target network provides stable target for learning. Without it, $y_t$ changes as $\theta$ changes, causing oscillations. By updating $\theta^-$ only periodically, we stabilize training.
\end{frame}

% ============================================================================
% SLIDE 14: DQN TRAINING LOOP
% ============================================================================
\begin{frame}{DQN Training Loop: 7 Steps per Experience}
  \scriptsize
  \begin{tabular}{|c|p{1.8cm}|p{4.4cm}|}
    \hline
    \textbf{Step} & \textbf{Action} & \textbf{Details} \\
    \hline
    1 & Observe state & AV measures $(d, c, r, w, v, m)$ \\
    \hline
    2 & Estimate Q & Feed state to network $\to$ get $Q(s, 0)$ and $Q(s, 1)$ \\
    \hline
    3 & Choose action & $\varepsilon$-greedy: explore with prob $\varepsilon$, exploit (max $Q$) with prob $1-\varepsilon$ \\
    \hline
    4 & Interact & Execute $a$; environment returns reward $r$ and next state $s'$ \\
    \hline
    5 & Store memory & Add $(s, a, r, s')$ to replay buffer (circular, $\sim 10k$) \\
    \hline
    6 & Sample/train & Mini-batch ($\sim 32$) from buffer; compute loss and update $\theta$ \\
    \hline
    7 & Repeat & Continue episodes until cumulative reward converges \\
    \hline
  \end{tabular}

  \vspace{0.08cm}
  \footnotesize\textbf{Why Replay Memory?} Random sampling decorrelates data, reducing variance and stabilizing learning.

  \normalsize
\end{frame}

% ============================================================================
% SLIDE 15: MDP SUMMARY
% ============================================================================
\begin{frame}{Summary: Complete MDP for iRDRC}
  \small
  \begin{tabular}{|p{1.5cm}|p{2.3cm}|p{1.7cm}|p{3.2cm}|}
    \hline
    \textbf{Component} & \textbf{Definition} & \textbf{Size/Form} & \textbf{Implementation Note} \\
    \hline
    State Space & $(d, c, r, w, v, m)$ & 352 discrete states & Fully observable MDP \\
    \hline
    Action Space & $\{0: \text{comm}, 1: \text{radar}\}$ & 2 actions per state & Binary choice each step \\
    \hline
    Reward & $a=0$: $+r_1=+2$ (c=0, $\ominus$), $+r_2=+1$ (c=1, $\ominus$), $-r_3=-50$ ($\oplus$); $a=1$: $+r_4(b+1)$ ($\oplus$), $0$ ($\ominus$) & $-50$ to $+25$ & Safety-first design \\
    \hline
    Transition & Queue dynamics deterministic; channel \& env factors stochastic & $P(\oplus)$ via Eq. 1 & Markov: indep. of history \\
    \hline
    Discount & $\gamma \approx 0.99$ (not specified) & Scalar in $[0,1)$ & Typical DRL value \\
    \hline
    Solver & DQN (MLP, 2 hidden layers) & Neural net weights & Converges in $\sim$170 episodes \\
    \hline
  \end{tabular}
\end{frame}

% ============================================================================
% SLIDE 16: KEY TAKEAWAYS
% ============================================================================
\begin{frame}{Key Takeaways: Why MDP + DQN Works}
  \begin{enumerate}
    \item \textbf{MDP Captures Structure:} Markov property holds; past doesn't matter, only current state. Future independent of history enables learning from experience.

    \item \textbf{Reward Design Encodes Priorities:} Rewards for \emph{successful} communication (\emph{no event}): $+1$ to $+2$. Penalty for \emph{missing hazard} in comm mode: $-50$. This 25--50× asymmetry enforces safety-first learning by hand-tuned weight.

    \item \textbf{Radar Reward Scales with Risk:} Detecting event while many factors unfavorable ($b=4$) earns $+25$; while safe ($b=0$) earns $+5$. Encourages radar mode exactly when it matters most.

    \item \textbf{Discount Factor Ensures Long-Horizon Planning:} $\gamma \approx 0.99$ makes policy value future safety. With $\gamma=0.5$, policy would comm recklessly (future 50\% discounted). With $\gamma=0.99$, policy balances.

    \item \textbf{DQN Scales Beyond Tabular:} 352 states OK for both Q-learning (280 episodes) and DQN (170 episodes). DQN faster via replay memory and mini-batch training.

    \item \textbf{Adaptation Emerges Automatically:} Policy not hand-coded. Learning discovers ``use radar when weather+speed risky, communicate when safe'' by optimizing cumulative reward.
  \end{enumerate}
\end{frame}

% ============================================================================
% SLIDE 17: MODELING GAPS
% ============================================================================
\begin{frame}{Modeling Gaps to Address Before Implementation}
  \small
  \begin{tabular}{|p{1.3cm}|p{2.0cm}|p{5.0cm}|}
    \hline
    \textbf{Gap} & \textbf{What's Missing} & \textbf{Action for Implementation} \\
    \hline
    Probability & $p^0_r, p^1_r, p^0_m, p^1_m$ not specified & Estimate: $p^0=0.01, p^1=0.05$ \\
    \hline
    State Evolution & $\tau_i$ (prob. favorable) not given & Assume $\tau_r \approx 0.7, \tau_w \approx 0.8, \tau_v \approx 0.6, \tau_m \approx 0.7$ \\
    \hline
    Discount $\gamma$ & Not specified in paper & Use $\gamma \approx 0.99$ (standard) \\
    \hline
    Episode Length $T$ & Not explicit; 500--1000 steps & Experiment with $T=500$, measure convergence \\
    \hline
    DQN Params & Learning rate, batch, buffer, $\varepsilon$-decay unspecified & Defaults: $\alpha=0.001, \text{batch}=32, \text{buffer}=10k$ \\
    \hline
    Network Arch & Hidden layers/units not specified & 2 layers, 64 units each $(6 \to 64 \to 64 \to 2)$ \\
    \hline
  \end{tabular}
\end{frame}

% ============================================================================
% FINAL SLIDE: THANK YOU
% ============================================================================
{
  \setbeamertemplate{background}{
    \begin{tikzpicture}[overlay, remember picture]
      \useasboundingbox (current page.south west) rectangle (current page.north east);
      \fill[white] (current page.south west) rectangle ([yshift=-0.9\beamer@height] current page.north east);
      \begin{scope}[x={(current page.north east)}, y={(current page.south west)}]
        \usebeamercolor{Feather}
        \fill[Feather.bg] ([yshift=-\beamer@height] current page.north west) rectangle (current page.north east);
      \end{scope}
    \end{tikzpicture}
  }
  \begin{frame}
    \vspace{2cm}
    \begin{center}
      \Huge\textbf{Thank You}

      \vspace{3cm}

      \Large\textit{Questions \& Discussion}
    \end{center}
  \end{frame}
}

\end{document}
